\section{Introduction}\label{sec:introduction}
\PARstart{V}{isual} tyre inspection combines several different questions. A system may locate a tyre, identify its tread region, estimate image-relative boundaries, classify a usage-related category, and expose the evidence behind its output. These functions have different supervision and different meanings. A good mask is evidence of region agreement, whereas a reliable physical-wear estimate requires suitable measurements and independent evaluation. Optical tread-depth systems such as TireEye explicitly connect image structure to a physical reference \cite{tireeye}. A classifier trained on mileage categories does not acquire that connection merely by producing a plausible heat map.

The present work grew from a compact multitask tyre-inspection system. Its dense outputs agreed with the annotations, but mileage-proxy classification transferred poorly to the allocated held-out tyres. Revisions addressed gradient interference, appearance sensitivity, reference exposure, spatial layout, and separation of geometry from texture learning. This sequence motivates a narrower question: \emph{which spatial summaries preserve the texture evidence that a fixed local decoder actually uses?}

This question matters because reconstruction and representation diagnostics can disagree for structural reasons. A local reconstruction head can read individual tokens without ever using a pooled image vector. Simultaneously, an auxiliary global head can be dominated by one direction of variation. Observing both properties does not constitute a contradiction. Nor does increasing a rank statistic automatically establish better evidence: coordinate scaling changes spectral balance, and a recorded mask layout can contribute variation unrelated to the desired texture content.

We study these issues in TRACE (Texture Reconstruction with Anchored Canonical Evidence), a small model built around frozen predicted geometry and fixed local image descriptors. Its six targets describe contrast, edge energy, orientation, curvature, and contrast occupancy. They are observable image properties, not annotated groove depth or damage. The model improves local and masked reconstruction against the declared controls but never satisfies its original representation-selection requirements. We preserve that outcome and examine the saved final tokens using the same frozen decoder, rather than retraining until a favorable endpoint appears.

The paper makes three contributions.
\begin{enumerate}
\item We formulate an equal-budget spatial audit for aligned texture tokens. A shared position template and mandatory constant mode yield a closed-form optimum for each train-only squared surrogate. This lets us compare input variance, decoder sensitivity, and decoded snapshots at the same number of input coefficients.
\item We derive a signed loss decomposition for exchanging spatial projection and the frozen decoder. Together with a weighted output-drift bound, it distinguishes the effect of spatial compression from changes introduced by the order of nonlinear decoding.
\item We test the audit on both TRACE reconstruction paths, across seven mode budgets, two recorded views, and perturbations to training-component composition. At five modes, learned spatial bases retain more reconstruction advantage than the tested fixed bases. Decoder-aware fitting offers no consistent gain over ordinary spatial PCA across the two held-out roles. The original interventions, covariance analysis, and mask-layout check explain why a global rank score cannot substitute for this measurement.
\end{enumerate}

The study uses 570 images of 12 physical tyres identified in the acquisition records. Each held-out role has two conservative components, and the historical test was inspected repeatedly during development. The evidence concerns a fixed TRACE model, specific descriptor targets, and two known allocations. Cross-tyre population inference and physical wear measurement require separate evidence.
